\documentclass[11pt]{article}
\usepackage[utf8]{inputenc}
\usepackage[T1]{fontenc}
\usepackage{amsmath,amssymb,amsthm}
\usepackage[margin=1in]{geometry}
\usepackage{hyperref}
\usepackage{booktabs}

\hypersetup{
  colorlinks=true,
  linkcolor=blue,
  citecolor=blue,
  urlcolor=blue
}

\title{\textbf{Formalizing the Reduced Mahler $3/2$ Dynamical System in Lean 4:\\
Exact 1D Slack Reduction, Constructive Coasting Survivors,\\
Conditional Finiteness Theorems, and Certified Transcendence Barriers}}

\author{Joshua Coon\\
\small (with Astra and Argon AI assistance via GitHub Copilot)}
\date{October 18, 2026\\
\small\texttt{https://github.com/jjc16/mahler-3-2-lean4}\\
\small Toolchain: Lean 4.19.0 + Mathlib v4.19.0 (114 compiled modules, 0 \texttt{sorry}s)}

\begin{document}
\maketitle

\begin{abstract}
Mahler's $3/2$ problem (1968) asks whether there exists any real ``$Z$-number'' $\alpha > 0$ such that the fractional part $\{(3/2)^n \alpha\}$ lies in $[0, 1/2)$ for every integer $n \ge 0$. Writing $\alpha = a_0 + \xi$ with $a_0 = \lfloor \alpha \rfloor \in \mathbb{N}$ and $\xi \in [0, 1)$, survival up to step $N$ is equivalent to the discrete Mahler orbit $g_{n+1} = \lfloor 3 g_n / 2 \rfloor$ ($g_0 = a_0$) admitting a non-empty intersection of $N$ nested archimedean companion intervals $[\mathrm{lowerBound}(a_0, n),\, \mathrm{upperBound}(a_0, n))$.

In this work, we present a 114-module, \texttt{sorry}-free formalization in Lean~4 that investigates both \textbf{constructive survival} and \textbf{obstructions to infinite orbits}:
\begin{enumerate}
  \item \textbf{Exact 1D Dynamical Reduction:} By grouping each odd step $g \mapsto (3g+1)/2$ with its $s(k)-1$ subsequent even steps $g \mapsto 3g/2$ via $3 b_k + 1 = 2^{s(k)} u(k)$ ($\gcd(u(k), 6) = 1$), we collapse all $\binom{K}{2}$ pairwise companion-interval inequalities into a \textbf{single 1D recurrence} (\texttt{minSlack}), and prove its exact two-way equivalence with the existence of a real companion $\xi \in \mathbb{R}$. Normalizing by $2^{T(m)}$ transforms the problem into an autonomous 1D piecewise-expanding interval map $\mu_2(s, m+1) = \min\big(1,\, (3/2)^{s(m)} \mu_2(s, m) - 2\big)$ on $(0, 1]$ governed by a disjoint L\"uroth--Markov partition $(0, 1] = \bigsqcup_{c=2}^\infty (2(2/3)^c,\, 3(2/3)^c] \cap (0, 1]$.
  \item \textbf{Constructive Survivors \& 113-Step Record-Holder:} We prove that any block sequence with $s(k) \ge 3$ saturates the ceiling $\mathrm{minSlack}\, s\, m = 2^{T(m)}$ and automatically admits a valid real companion. Using this sufficiency theorem, we construct closed-form seed families surviving arbitrarily many blocks prior to ceiling exit, and machine-certify a 55-bit integer ($\mathrm{recordHolderA0} = 29{,}871{,}555{,}608{,}238{,}677 < 2^{55}$) that survives \textbf{113 steps} (58 steps past its 55-bit binary length) before dying at step 114.
  \item \textbf{Unconditional Finiteness in Restricted Regimes:} We prove explicit, logarithm-free linear step cutoffs ruling out infinite orbits in four major regimes: (i) bounded-block runs ($s(k) \le c$), (ii) periodic, eventually periodic, or low-subword-complexity block sequences ($\mathrm{blockSteps}^{(r_1)} L \le T(r_2) + W$), (iii) quotients $u(m)$ supported on a fixed finite prime set ($S$-units, via Baker--Matveev), and (iv) linearly growing blocks ($s(m) \ge \delta T(m)$) or exponential slack pinches ($\Delta_m \le 2^{(1-\delta)T(m)}$, via Ridout).
  \item \textbf{Certified Diagnosis of Global Transcendence Barriers (Dead Ends 1--18):} We formally prove \emph{why} standard elementary, polynomial-determinant, Baker--Matveev, Ridout, Schlickewei Subspace, Stewart--Bugeaud base-digit, and Pad\'e--Hankel methods stall on a general aperiodic sequence $s$: each method hits an exact algebraic conservation law (rank-1 determinant factorization, linear Weil height growth $h(u(m+1)/u(m)) \sim T(m)\log(3/2)$, exact tangency at Ridout's critical $Q_m^{-2}$ exponent via $(1 - \frac{\log 2}{\log 3}) + \frac{\log 2}{\log 3} + 1 = 2$, or a rank-2 hyperplane syzygy across shifted origins).
\end{enumerate}
\end{abstract}

\section{Introduction and Literature Attribution}

When attempting to prove Mahler's $3/2$ conjecture or related Collatz-type problems, informal proof sketches typically fall into two traps: (1) treating 2-adic parity constraints and archimedean interval constraints as statistically independent, or (2) constructing auxiliary linear combinations, Wronskians, or multi-origin determinants without noticing that the underlying 1D affine recurrence forces those combinations to be algebraically dependent. By formalizing every proposed attack vector in Lean~4, we separated genuine number-theoretic structure from mirages.

\subsection{Attribution to Prior Mathematical Literature}
\begin{enumerate}
  \item \textbf{Mahler (1968) \cite{Mah68}:} Introduced $Z$-numbers and proved that for each integer $a_0 \ge 1$, there is at most one real companion $\xi \in [0, 1)$ such that $\alpha = a_0 + \xi$ is a $Z$-number, governed by $g_{n+1} = \lfloor 3 g_n / 2 \rfloor$ and nested archimedean intervals (\texttt{MahlerDemandA.lean}, \texttt{PureIntegerPinch.lean}). Mahler also derived the series representation of $\xi_\infty$ and applied Ridout's theorem \cite{Rid57} to two-seed differences (\texttt{RidoutSubspaceFramework.lean}).
  \item \textbf{Vijayaraghavan (1940), Pollington (1979), Choquet (1980) \& Flatto--Lagarias--Pollington (1995) \cite{Vij40,Pol79,Cho80,FLP95}:} Proved the $1/3$ fractional-part range theorem $\limsup_{n\to\infty} \{(3/2)^n \alpha\} - \liminf_{n\to\infty} \{(3/2)^n \alpha\} \ge 1/3$, which underpins our $1/3$ Archimedean Drain Budget ($\sum_{k=1}^m (2/3)^{T(k)+1} < 1/3$ in \texttt{TelescopingSeriesDuality.lean}) and explains why $s(k) = 1$ is forbidden (\texttt{EscapeBranchTwoCap.lean}). Flatto, Lagarias, and Pollington also established the symbolic-dynamical forbidden-subword hierarchy and average step-rate bounds that we rediscovered and machine-checked in odd-block length coordinates (\texttt{GlobalStepRateBounds.lean}, \texttt{MinSlackDynamicalSystem.lean}).
  \item \textbf{Classical Diophantine \& Transcendence Machinery \cite{Rid57,Sch77,Ste80,BW93,Mat00,CZ02,CZ03,AB07,Dub06}:} Adapted in Modules 89--102 to analyze linear forms in logarithms, multi-place Ridout products, Schlickewei Subspace systems, simultaneous base-$2$/base-$3$ representations, and multi-step GCDs.
\end{enumerate}

\subsection{Comparison with Prior Lean Formalizations \& Novel Contributions}
\begin{itemize}
  \item \textbf{Prior Lean Repositories:} Google DeepMind's \texttt{formal-conjectures} repository \cite{FC26} (\texttt{Mahler32.lean}, \texttt{Problem10\_9.lean}, \texttt{ZNumber.lean}) previously formalized the definition \texttt{IsZNumber}, the 5-line sanity check $\neg\mathrm{IsZNumber}\,1$, and \texttt{sorry}-ed problem statements for Mahler's conjecture and $\Omega(p/q) \ge 1/p$. Other Lean repositories bearing Mahler's name address distinct theorems (\texttt{giuliocaflisch/mahler} on $p$-adic continuous functions, \texttt{DiegoMarquesMath/MahlerLean} on Liouville numbers). Our 114-module library is the first proved Lean~4 development of the discrete Mahler $3/2$ dynamical system and its symbolic/Diophantine structure.
  \item \textbf{Reduced Coprime Quotient System \& 1D \texttt{minSlack} Collapse:} Formulating the dynamics via $2^{s(k+1)} u(k+1) = 3^{s(k)} u(k) + 1$ ($\gcd(u(k), 6) = 1$) with 3-adic block-boundary factorization $b_{k+1} = 3^{s(k)-1} u(k)$, collapsing all $\binom{K}{2}$ shifted companion inequalities into the 1D \texttt{minSlack} recurrence (\texttt{ConverseCompanionReconstruction.lean}), and normalizing it to the autonomous map $\mu_2(s, m+1) = \min\big(1,\, (3/2)^{s(m)} \mu_2(s, m) - 2\big)$ on the disjoint L\"uroth--Markov partition $(0, 1] = \bigsqcup_{c=2}^\infty (2(2/3)^c, 3(2/3)^c] \cap (0, 1]$ (\texttt{CoupledFractionalPartDynamics.lean}).
  \item \textbf{Additive Separation of Shifted Slacks \& Three-Origin Subspace Syzygy:} Proving $\frac{2^{T(r)} \Delta^{(r)}_{m-r}}{3^{T(m)}} = \frac{\Delta^{(0)}_m}{3^{T(m)}} - \frac{\Delta^{(0)}_r - 2^{T(r)}}{3^{T(r)}}$, eliminating \texttt{shiftSeq} from $0 < \mathrm{minSlack}\, s\, m$, and proving the 3-origin Schlickewei hyperplane syzygy (\texttt{MultiOriginSubspaceSyzygy.lean}).
  \item \textbf{Carryless Base-3 Frozen-Prefix Equivalence \& Dual Remainder Law:} Proving $0 < \Delta_m \iff \lfloor 2^{T(m+1)+1} u(m) / 3^{T(m)} \rfloor = 6a_0 + 2$ and $(P_m \bmod 3^{T(m)}) + (Q_m \bmod 2^{T(m)}) = 3^{T(m)}$ (\texttt{SimultaneousBaseDigitRigidity.lean}).
  \item \textbf{Three-Way Multi-Step GCD Identities \& Two-Sided CRT Floors:} Proving $\gcd(u(r), u(r+d)) = \gcd(u(r), R^{(r)}_d)$, $\gcd(\Delta_r, \Delta_{r+d}) = \gcd(\Delta_r, R^{(r)}_d)$ (with $\gcd(\Delta_r, \Delta_{r+1}) = 1$), $\gcd(u(m), \Delta_m) = \gcd(u(m), 2a_0+1) = \gcd(\Delta_m, 2a_0+1)$, and the two-sided CRT floors $u(k) \ne 1$ ($k \ge 1$), $1735 \le u(k)$ at $s(k) = 2$, $65 \le u(k)$ across $2 \le s(k) \le 6$, and $103 \le u(k)$ on ceiling-saturated runs (\texttt{MultiStepQuotientGCD.lean}, \texttt{ThreeAdicAndSeedGCDLocks.lean}, \texttt{Mod2304CombinedSieve.lean}, \texttt{TwoSidedCRTAndBackwardObstructions.lean}, \texttt{Mod216TwoSidedQuotientFloor.lean}).
  \item \textbf{Discovery \& Certification of the 55-Bit, 113-Step Coasting Survivor:} Constructing and verifying $\mathrm{recordHolderA0} = 29{,}871{,}555{,}608{,}238{,}677 < 2^{55}$, which survives 113 steps and dies at step 114 (\texttt{MahlerGrandSynthesis.lean}).
\end{itemize}

\section{The Odd-Block Reduction \& Dual 2-Adic / 3-Adic Recurrence}

A sequence of block lengths $s : \mathbb{N} \to \mathbb{N}$ and odd quotients $u : \mathbb{N} \to \mathbb{N}$ forms a \texttt{ValidBlockChain} $a_0\, s\, u\, K$ if for every $k < K$:
\[
  1 \le s(k), \qquad u(k) \bmod 2 = 1, \qquad 3 \cdot b_k + 1 = 2^{s(k)} u(k),
\]
where $b_k := \mathrm{orbit}\, a_0\, (T(k))$ and $T(k) := \mathrm{blockSteps}\, s\, k = \sum_{i=0}^{k-1} s(i)$. Across consecutive blocks ($k + 1 < K$), the quotients and orbit boundaries satisfy:
\[
  2^{s(k+1)} u(k+1) = 3^{s(k)} u(k) + 1, \qquad b_{k+1} = 3^{s(k)-1} u(k), \qquad \gcd(u(k), 6) = 1.
\]
Defining $R(s, 0) = 0$ and $R(s, m+1) = 3^{s(m)} R(s, m) + 2^{\mathrm{tailSteps}\, s\, m}$, the intrinsic integer slack satisfies:
\[
  \Delta_m := \mathrm{integerSlack}\, s\, m := 3^{T(m)} - 2^{s(0)+1} R(s, m) = 3^{T(m)+1}(2a_0 + 1) - 2^{T(m+1)+1} u(m).
\]

\section{Exact 1D Dynamical Reduction \& Converse Companion Reconstruction}

Defining $\mathrm{minSlack}\, s\, 0 = 1$ and $\mathrm{minSlack}\, s\, (m+1) = \min\big(2^{T(m+1)},\, 3^{s(m)}\mathrm{minSlack}\, s\, m - 2^{T(m+1)+1}\big)$, we proved in \texttt{ConverseCompanionReconstruction.lean}, \texttt{MultiOriginSubspaceSyzygy.lean}, \texttt{SimultaneousBaseDigitRigidity.lean}, and \texttt{CoupledFractionalPartDynamics.lean} the four-way equivalence:
\begin{align*}
  \big(\exists \xi \in \mathbb{R},\ \mathrm{ValidCompanionUpTo}\, a_0\, \xi\, (T(K))\big)
  &\iff 0 < \mathrm{minSlack}\, s\, (K - 1) \\
  &\iff \forall j \le K - 1,\ 3^{T(K-1)}\big(\Delta_j - 2^{T(j)}\big) < 3^{T(j)} \Delta_{K-1} \\
  &\iff \forall j \le K - 1,\ \left\lfloor \frac{2^{T(K) - T(j) + 1} u(K - 1)}{3^{T(K-1) - T(j)}} \right\rfloor = 6 b_j + 2 \\
  &\iff \mu_2(s, K - 1) := \frac{\mathrm{minSlack}\, s\, (K-1)}{2^{T(K-1)}} \in (0, 1],
\end{align*}
where $\mu_2(s, m+1) = \min\big(1,\, (3/2)^{s(m)} \mu_2(s, m) - 2\big)$ evolves on the pairwise-disjoint L\"uroth--Markov partition $(0, 1] = \bigsqcup_{c=2}^\infty (2(2/3)^c,\, 3(2/3)^c] \cap (0, 1]$.

\section{Certified Transcendence \& Diophantine Barriers (Dead Ends 1--18)}

\begin{enumerate}
  \item \textbf{Dual-Place Convergence in $\mathbb{Z}_2 \times \mathbb{R}$ (\texttt{MahlerSeriesHankelFramework.lean}):} The generating series $\Sigma_m(s) = \sum_{k=1}^m (2/3)^{T(k)+1} = \frac{2^{s(0)+1} R(s, m)}{3^{T(m)+1}}$ satisfies $\Sigma_m(s) + \frac{6a_0+2}{3} = \frac{2^{T(m+1)+1} u(m)}{3^{T(m)+1}}$, converging to $-\frac{6a_0+2}{3}$ in $\mathbb{Q}_2$ while staying $< 1/3$ in $\mathbb{R}$.
  \item \textbf{Determinant \& Pad\'e--Hankel Collapse (Dead Ends 12 \& 18, \texttt{DeterminantFactorization.lean}, \texttt{MahlerSeriesHankelFramework.lean}):}
  \[
    3^{T(r)+1} P_{r+d} - 3^{T(r+d)+1} P_r = 3^{T(r)+1} \cdot 2^{T(r+1)+1} \cdot R(\mathrm{shiftSeq}\, s\, r,\, d),
  \]
  dropping the 2-adic valuation from $T(r+d+1)+1$ back to $T(r+1)+1$.
  \item \textbf{Baker--Matveev \& Multi-Step GCD Weil Height Growth (Dead Ends 13 \& 17, \texttt{BakerMatveevFramework.lean}, \texttt{MultiStepQuotientGCD.lean}):} Because $\gcd(u(r), u(r+d)) = \gcd(u(r), R^{(r)}_d) < 3^{S_{r, d}} / 2^{s(r)+1}$, the reduced denominator satisfies $3^{T(r)}(3a_0 + 1) < 2^{T(r)} 3^{S_{r, d}} (u(r)/\gcd(u(r), u(r+d)))$, forcing linear Weil height growth $h \sim T(r)\log(3/2)$.
  \item \textbf{Ridout Critical $Q_m^{-2}$ Tangency (Dead End 14, \texttt{RidoutSubspaceFramework.lean}):} In any sub-ceiling phase, $2^{T(m)}$ cancels identically from $u(m) \cdot \Delta_m$ ($2^{s(m)+D+1} u(m) \Delta_m < Q_m$), placing the 3-place $\{\infty, 2, 3\}$ Ridout product at exponent $(1 - \frac{\log 2}{\log 3}) + \frac{\log 2}{\log 3} + 1 = 2$.
  \item \textbf{Schlickewei Subspace Hyperplane Syzygy (Dead End 15, \texttt{MultiOriginSubspaceSyzygy.lean}):} Any three shifted origins $(0, r_1, r_2)$ satisfy the $m$-independent hyperplane syzygy
  \[
    D_0(r_2) \cdot 3^{T(r_1)} \big(2^{T(r_1)} \Delta^{(r_1)}_{m - r_1} - \Delta^{(0)}_m\big) = D_0(r_1) \cdot 3^{T(r_2)} \big(2^{T(r_2)} \Delta^{(r_2)}_{m - r_2} - \Delta^{(0)}_m\big).
  \]
\end{enumerate}

\section*{Acknowledgments}
This research and its 114-module Lean~4 formalization were developed by Joshua Coon with the assistance of the \textbf{Astra} and \textbf{Argon} AI models in GitHub Copilot.

\begin{thebibliography}{99}
\bibitem{Mah68} K. Mahler, ``An unsolved problem on the powers of $3/2$,'' \emph{J. Austral. Math. Soc.} \textbf{8} (1968), 313--321.
\bibitem{Rid57} D. Ridout, ``Rational approximations to algebraic numbers,'' \emph{Mathematika} \textbf{4} (1957), 125--131.
\bibitem{Vij40} T. Vijayaraghavan, ``On the fractional parts of the powers of a number (I),'' \emph{J. London Math. Soc.} \textbf{15} (1940), 159--160.
\bibitem{Pol79} A. D. Pollington, ``On the density of sequence $\{\theta^k\}$,'' \emph{Illinois J. Math.} \textbf{23} (1979), 511--515.
\bibitem{Cho80} G. Choquet, ``$\theta$-ferm\'es et dimension de Hausdorff,'' \emph{C. R. Acad. Sci. Paris S\'er. A--B} \textbf{290} (1980), A863--A868.
\bibitem{FLP95} L. Flatto, J. C. Lagarias, and A. D. Pollington, ``On the range of fractional parts $\{\xi (p/q)^n\}$,'' \emph{Acta Arith.} \textbf{70} (1995), 197--227.
\bibitem{Sch77} H. P. Schlickewei, ``The $\mathfrak{p}$-adic Thue--Siegel--Roth--Schmidt theorem,'' \emph{Arch. Math.} \textbf{29} (1977), 267--270.
\bibitem{Ste80} C. L. Stewart, ``On the representation of an integer in two different bases,'' \emph{J. Reine Angew. Math.} \textbf{319} (1980), 63--72.
\bibitem{BW93} A. Baker and G. W\"ustholz, ``Logarithmic forms and group varieties,'' \emph{J. Reine Angew. Math.} \textbf{442} (1993), 19--62.
\bibitem{Mat00} E. M. Matveev, ``An explicit lower bound for a homogeneous rational linear form in the logarithms of algebraic numbers. II,'' \emph{Izv. Math.} \textbf{64} (2000), 1217--1269.
\bibitem{CZ02} P. Corvaja and U. Zannier, ``Some new applications of the Subspace Theorem,'' \emph{Compositio Math.} \textbf{131} (2002), 319--340.
\bibitem{CZ03} P. Corvaja and U. Zannier, ``On the greatest prime factor of $(ab+1)(ac+1)$,'' \emph{Proc. Amer. Math. Soc.} \textbf{131} (2003), 1705--1709.
\bibitem{AB07} B. Adamczewski and Y. Bugeaud, ``On the complexity of algebraic numbers I. Expansions in integer bases,'' \emph{Ann. of Math.} \textbf{165} (2007), 547--565.
\bibitem{Dub06} A. Dubickas, ``On the distance from a rational power to the nearest integer,'' \emph{J. Number Theory} \textbf{117} (2006), 222--239.
\bibitem{FC26} The Formal Conjectures Authors (Google DeepMind), ``Formal Conjectures in Lean 4: \texttt{Mahler32.lean}, \texttt{Problem10\_9.lean}, and \texttt{ZNumber.lean},'' \url{https://github.com/google-deepmind/formal-conjectures} (2025--2026).
\end{thebibliography}

\end{document}
